% !TEX root = ../main.tex
\section*{Introducción}

[Una cuartilla: 1) campo de investigación, 2) lo ya estudiado por otros,
3) área de oportunidad, 4) objetivo del proyecto, 5) justificación.]

\newpage

\section{Planteamiento del problema}
Una organización del sector público mantiene su sistema principal, el cual sostiene gran parte de la
operación administrativa y operativa e integra las diferentes áreas de esta organización. La información
se almacena en dos motores de base de datos, uno relacional y otro documental. El segundo, al no ser relacional,
no implementa llaves foráneas como restricción del propio motor y, al tratarse de motores distintos, tampoco
hay un mecanismo de integridad referencial que los conecte entre sí, por lo que la relación entre los datos
depende de una convención y no de una restricción verificada por el sistema.

El sistema nació hace tres años orientado principalmente a elementos operativos, con funciones 
como la generación de resguardos de unidad y bitácoras de servicio. Con el tiempo escaló hacia el área
administrativa de la organización, incorporando de forma progresiva nuevas funcionalidades hasta
convertirse en el sistema principal de la organización, sin que su arquitectura se revisara al mismo ritmo
que su alcance. Actualmente sostiene operaciones críticas de la organización, mientras que su
proceso de despliegue continúa siendo completamente manual.

Para realizar un despliegue es necesario acceder mediante SSH al servidor donde se aloja la
aplicación, actualizar el código mediante operaciones de Git sobre el repositorio, ejecutar
manualmente las migraciones pendientes y limpiar cachés de configuración y rutas, sin mecanismos
de verificación automática. El historial del repositorio refleja un volumen considerable de
cambios: desde su inicio se han fusionado 563 solicitudes de cambio en la rama principal, sobre
aproximadamente 8,000 commits. Estas integraciones pueden incluir cambios que deben aplicarse
sobre los motores de base de datos: modificación de esquemas en el motor relacional y, en el
motor no relacional, que no tiene esquema fijo, operaciones equivalentes como la creación de
índices; en ambos casos deben identificarse y ejecutarse manualmente y en el orden correspondiente.

El ritmo de integración de cambios a la rama principal no equivale a la frecuencia real de
despliegue: con regularidad se agrupan varios cambios fusionados antes de llevarlos al servidor,
aunque no siempre ocurre así ni existe un criterio fijo de cuántos se acumulan. Tampoco existe un
registro sistemático del tiempo entre la fusión de un cambio a la rama principal y su despliegue
real en producción; de forma informal se sabe que ese tiempo es corto, ya que no se deja código
fusionado sin desplegar por periodos prolongados, pero no hay un dato preciso que lo respalde.
Tampoco se conoce, con el mismo nivel de precisión, qué proporción de despliegues requiere una
corrección posterior.

El proceso manual presenta dos modos de falla asociados a la ejecución de pasos a mano. El
primero, ya observado de forma recurrente, es olvidar limpiar las cachés de configuración y rutas
al final de cada despliegue, lo que deja el sistema sirviendo una versión desactualizada hasta que
se detecta el problema. El segundo, un riesgo identificado pero sin evidencia de haberse
materializado, es específico de los despliegues que agrupan varios cambios: omitir alguna
migración o seeder de los distintos cambios acumulados.

\subsection{Antecedentes}
El primer elemento a considerar al intervenir un proceso de entrega de software es establecer un mecanismo
sistemático y validado para medir su desempeño antes y después de la automatización. Forsgren, Humble y Kim
\cite{Forsgren2018} sintetizaron años de investigación empírica en miles de organizaciones y aislaron cuatro 
métricas clave, conocidas en la literatura como métricas DORA: frecuencia de despliegue (\emph{deployment frequency}), 
tiempo de entrega de los cambios (\emph{lead time for changes}), tiempo de restauración del servicio 
(\emph{time to restore service}) y tasa de fallos por cambio (\emph{change failure rate}). El hallazgo
central demuestra que las organizaciones de alto desempeño no sacrifican estabilidad por velocidad, sino
que consiguen ambas de forma simultánea. Para este trabajo, estas métricas proporcionan un vocabulario y una
definición operacional ya validados, ofreciendo un marco científico para cuantificar la mejora respecto a la línea base
(\emph{baseline}) de los despliegues manuales.

Sobre la adopción de Integración y Entrega Continua (CI/CD), existe un consenso sólido respecto a las principales barreras
de implementación. Revisiones sistemáticas de literatura empírica desarrolladas por Lukkanen et al. \cite{Laukkanen2017} y
Proulx et al. \cite{Proulx2018} coinciden en que los factores humanos y organizacionales superan holgadamente a los
problemas técnicos en la generación de fallas durante la liberación de código. Estos hallazgos respaldan directamente
el diagnóstico del sistema de estudio de este proyecto: la omisión de limpieza de cachés tras un despliegue
manual constituye una falla observada de forma recurrente, mientras que la omisión de migraciones al agrupar
múltiples \emph{commits} o cambios constituye un riesgo potencial identificado, sin evidencia de haberse
materializado. Ambos problemas son consistentes con la influencia de factores humanos. La literatura señala
que la automatización mediante un \emph{pipeline} establece un mecanismo restrictivo de construcción que puede
reducir la posibilidad de errores humanos por omisión al hacer explícitos y verificables los pasos del despliegue.

Al intentar aplicar estos principios a sistemas \emph{legacy}, la literatura técnica se inclina hacia una refactorización
profunda o la descomposición en microservicios. Sin embargo, diferentes estudios recientes en la ingeniería
de software empresarial advierten sobre los elevados riesgos y costos que puede ocasionar esta estrategia,
especialmente en infraestructura restringida y autoalojada. La práctica documentada respalda la modernización
incremental \cite{Kosinska2020}, un enfoque que prioriza la contenedorización de la aplicación tal como está,
aislando sus dependencias de tiempo de ejecución (como el \emph{runtime} de la aplicación) sin reescribir la lógica
del negocio. Este es un paso crítico en contextos institucionales que operan sobre infraestructura propia y carecen de
orquestación en la nube pública.

Un desafío arquitectónico particular del sistema en cuestión es su dependencia de un esquema de base de datos dual 
y heterogéneo (un motor relacional y uno documental simultáneos), desprovisto de un mecanismo nativo de integridad 
referencial. Estudios empíricos sobre entrega continua, como el de Shahin et al. \cite{Shahin2017}, identifican la 
sincronización y la evolución de esquemas de bases de datos como uno de los cuellos de botella técnicos más severos
al automatizar despliegues. Dado que no existe una gestión de transacciones distribuidas entre ambos motores, 
cualquier estrategia de automatización debe tratar ambas bases como una unidad lógica unificada al ejecutar 
migraciones o \emph{backfills} de datos. Esto impone una restricción de diseño severa para el \emph{pipeline} propuesto.

Por otra parte, el obstáculo más significativo frente a la literatura tradicional de CI/CD es la 
validación del código. La inmensa mayoría de los trabajos académicos asumen la preexistencia de una 
suite de pruebas automatizadas como el \emph{quality gate} (compuerta de calidad) principal. 
Cuando un sistema \emph{legacy} carece de pruebas unitarias o de integración, forzar la entrega continua 
presenta un riesgo crítico. Para cerrar esta brecha en escenarios con datos altamente acoplados, 
la industria sugiere la instanciación de entornos de prueba bajo demanda 
(\emph{on-demand preview environments}), nutridos con \emph{dumps} sanitizados de producción. 
Esto transfiere la responsabilidad de la validación del \emph{pipeline} a un entorno aislado donde 
se puede realizar una comprobación funcional manual de alta fidelidad.

El contexto operativo institucional del sector público con servidores propios (\emph{on-premise})
añade una capa final de complejidad debido a restricciones presupuestales y la falta de registros de 
contenedores (\emph{container registries}) preexistentes. Un precedente altamente relevante en este rubro 
es el trabajo de Morales Agredo \cite{MoralesAgredo2021} en la Pontificia Universidad Javeriana Cali, quien 
diseñó una propuesta de CI/CD orientada a aplicaciones \emph{legacy} universitarias. Su investigación 
demuestra que la transición hacia flujos automatizados con plataformas autoalojadas (\emph{self-hosted}) 
es viable e imperativa en organizaciones con infraestructura limitada, sin depender de proveedores externos.

Cruzando los ejes anteriores emerge con claridad la brecha técnica y operativa que motiva este trabajo. 
La literatura revisada documenta sólidamente cómo evaluar el desempeño, los riesgos del error humano, 
la viabilidad de la contenedorización incremental y la necesidad de modernización en instituciones 
públicas. Ninguno de estos trabajos, sin embargo, aborda de manera conjunta la automatización de la
entrega de un sistema monolítico con un esquema de datos heterogéneo, operado sobre infraestructura
\emph{on-premise} sin acceso a nube pública, y que deba garantizar la seguridad del despliegue en 
ausencia total de una suite de pruebas automatizadas. Es precisamente en la resolución simultánea
de estas restricciones operativas donde se inscribe la contribución de este proyecto.


\subsection{Definición del problema}
El problema central que aborda este trabajo radica en la inestabilidad y el riesgo técnico generados por la 
gestión manual del ciclo de liberación de software en un entorno monolítico en producción. La ejecución directa 
de rutinas de despliegue y configuración sobre la infraestructura en vivo, sumada a la imposibilidad de validar 
los cambios en entornos aislados de forma previa, expone la operación a fallas recurrentes inducidas por 
el factor humano y compromete la confiabilidad técnica de cada integración.

\subsection{Justificación}
La automatización del ciclo de entrega de software no es únicamente una medida para incrementar la velocidad en
el desarrollo, sino un mecanismo que permite mitigar riesgos operativos. En el contexto de este proyecto, la
ejecución de despliegues manuales sobre una infraestructura \emph{on-premise} expone a la organización
a posibles interrupciones del servicio derivadas de errores en la configuración, como la falta de sincronización
al aplicar migraciones o la simple omisión en la limpieza de cachés. Al carecer de una suite de pruebas
automáticas, el riesgo de integrar cambios defectuosos a producción es elevado.

Resolver este problema genera un impacto operativo directo: al sustituir las intervenciones manuales por un 
\emph{pipeline} declarativo y generar entornos de prueba bajo demanda, se crea una ``compuerta de calidad''
(\emph{quality gate}) que aísla los fallos antes de que alcancen a los usuarios finales. Esto beneficia
directamente al equipo técnico, al reducir la dependencia de procedimientos manuales susceptibles 
a errores de omisión o ejecución, y al favorecer la estandarización, reproducibilidad y trazabilidad 
del proceso de liberación. Asimismo, beneficia a la institución, al disminuir el riesgo de interrupciones
durante los despliegues y contribuir a la continuidad operativa de un sistema que sustenta su gestión 
administrativa.


\subsection{Objetivo general}
Diseñar y validar empíricamente una arquitectura de integración y entrega continua (CI/CD) adaptada a sistemas
\emph{legacy}, mediante contenedorización y un entorno de prueba bajo demanda, con el propósito de mitigar la vulnerabilidad
operativa y reducir las intervenciones manuales rutinarias en la infraestructura de producción de un monolito operado 
sobre un esquema de base de datos dual y heterogéneo.

\subsection{Objetivos específicos}
Para alcanzar el objetivo general propuesto, se plantea estructurar las siguientes metas parciales:
\begin{itemize}
  \item Aislar el entorno de ejecución de la aplicación mediante la adopción de imágenes de contenedor inmutables, preservando
  la lógica del negocio del monolito sin recurrir a una reescritura de código orientada a microservicios.
  \item Orquestar un mecanismo de aprovisionamiento declarativo de entornos de prueba efímeros bajo demanda,
  alimentados por volcados de datos sanitizados del esquema heterogéneo (MySQL y MongoDB) de producción, para proveer un
  entorno de validación de alta fidelidad.
  \item Diseñar e implementar un \emph{pipeline} de Integración Continua (CI) sobre infraestructura autoalojada, 
  que incorpore compuertas de calidad (\emph{quality gates}) basadas en análisis estático avanzado, para verificar
  automáticamente cada cambio integrado al repositorio y detectar defectos, vulnerabilidades y \emph{code smells} 
  antes de su despliegue, ante la ausencia histórica de una suite de pruebas automatizadas.
  \item Evaluar comparativamente la resiliencia operativa de la arquitectura propuesta frente a la línea base de despliegues
  manuales, cuantificando la mejora a través de métricas estandarizadas (\emph{deployment frequency} y \emph{change failure rate}).
\end{itemize}

\subsection{Hipótesis}
[Opcional según el tipo de trabajo.]

\subsection{Alcances y limitaciones}
[Qué sí y qué no cubre el proyecto.]
